\cleardoublepage

\setstretch{1.15}

\addcontentsline{toc}{chapter}{Abstract in Greek}
\chapter*{\selg{} {\fontfamily{udidot}\selectfont Περίληψη}}

 {\fontfamily{udidot}\selectfont \selg

  Οι εφαρμογές που επεξεργάζονται μεγάλους όγκους δεδομένων στην εικονική μηχανή \selam Java (JVM)\selg{} βασίζονται σε μεγάλες κρυφές μνήμες αποτελεσμάτων για να αποφεύγουν την επανάληψη υπολογισμών και στην κρυφή μνήμη σελίδων του λειτουργικού συστήματος (\selam I/O cache\selg) για να περιορίζουν τις λειτουργίες εισόδου/εξόδου προς τα αποθηκευτικά μέσα.
  Καθώς αυξάνεται ο όγκος των δεδομένων, οι εφαρμογές συχνά χρειάζονται μεγαλύτερες κρυφές μνήμες αποτελεσμάτων και μεγαλύτερους σωρούς στη \selam JVM\selg.
  Η αύξηση αυτή συνεπάγεται υψηλό κόστος, καθώς (1) η μνήμη \selam DRAM\selg{} είναι ακριβή και η επέκταση της χωρητικότητάς της είναι δύσκολη και (2) τα αντικείμενα που παραμένουν για μεγάλο χρονικό διάστημα στην κρυφή μνήμη μπορούν να αυξήσουν το κόστος συλλογής απορριμμάτων (\selam GC\selg).
  Για να περιορίσουν το κόστος αυτό, οι εφαρμογές συχνά μεταφέρουν τις κρυφές μνήμες αποτελεσμάτων εκτός του σωρού της \selam JVM\selg{} (\selam off-heap\selg), σε επίπεδα αποθήκευσης που διαχειρίζονται οι ίδιες και τα οποία χρησιμοποιούν γρήγορα μέσα, όπως οι δίσκοι \selam NVMe SSD\selg, με μεγαλύτερη χωρητικότητα και χαμηλότερο κόστος ανά \selam GB\selg{} από τη \selam DRAM\selg.
  Ωστόσο, η πρόσβαση σε αυτές τις κρυφές μνήμες συνήθως απαιτεί σειριοποίηση και αποσειριοποίηση (\selam S/D\selg), επιβαρύνοντας την επεξεργασία.
  Παράλληλα, ο σωρός της \selam JVM\selg{} και η κρυφή μνήμη σελίδων ανταγωνίζονται για τη διαθέσιμη \selam DRAM\selg: η παραχώρηση περισσότερης μνήμης στον σωρό μπορεί να μειώσει το κόστος συλλογής απορριμμάτων, αλλά αφήνει λιγότερη μνήμη για την προσωρινή αποθήκευση σελίδων από τα αποθηκευτικά μέσα.
  Καθώς οι εφαρμογές μεταβαίνουν μεταξύ φάσεων εκτέλεσης με διαφορετικές απαιτήσεις μνήμης και εισόδου/εξόδου, μια στατική κατανομή της \selam DRAM\selg{} δεν μπορεί να εξισορροπήσει αποτελεσματικά τα δύο αυτά κόστη.

  Η παρούσα διατριβή διερευνά πώς μπορεί να επεκταθεί ο σωρός της \selam JVM\selg{} σε αποθηκευτικά μέσα χωρίς δαπανηρές λειτουργίες εισόδου/εξόδου λόγω συλλογής απορριμμάτων, καθώς και πώς μπορεί να κατανεμηθεί δυναμικά η \selam DRAM\selg{} μεταξύ του σωρού και της κρυφής μνήμης σελίδων του λειτουργικού συστήματος.
  Αντιμετωπίζουμε τρεις συμπληρωματικές προκλήσεις: τη διαχείριση των αντικειμένων στο αποθηκευτικό μέσο χωρίς σάρωση ολόκληρου του σωρού που βρίσκεται σε αυτό, την ενσωμάτωση της μεταφοράς αντικειμένων στη συλλογή απορριμμάτων με σύντομες παύσεις και τη δυναμική προσαρμογή της κατανομής της \selam DRAM\selg{} μεταξύ του σωρού της \selam JVM\selg{} και της κρυφής μνήμης σελίδων.

  Αρχικά, σχεδιάζουμε έναν σωρό σε αποθηκευτικό μέσο για να βελτιώσουμε την απόδοση εφαρμογών που δίνουν προτεραιότητα στη ρυθμαπόδοση.
  Επεκτείνουμε έναν συλλέκτη απορριμμάτων που στοχεύει στην υψηλή ρυθμαπόδοση, ώστε να υποστηρίζει μια αρχιτεκτονική δύο σωρών: (1) ενός σωρού στη \selam DRAM\selg{} (\selam DH\selg) και (2) ενός σωρού σε \selam NVMe SSD\selg{} (\selam SH\selg), οργανωμένου σε περιοχές και αντιστοιχισμένου στον χώρο διευθύνσεων της \selam JVM\selg.
  Η αρχιτεκτονική παρέχει στις εφαρμογές έναν ενιαίο, διαχειριζόμενο σωρό, αποκρύπτοντας τα επιμέρους επίπεδα μνήμης, ενώ τα αντικείμενα στον \selam SH\selg{} παραμένουν άμεσα προσπελάσιμα χωρίς σειριοποίηση και αποσειριοποίηση.
  Παρέχουμε μια διεπαφή υποδείξεων, μέσω της οποίας οι εφαρμογές προσδιορίζουν αντικείμενα προς μεταφορά στον \selam SH\selg, χωρίς να διαχειρίζονται οι ίδιες την τοποθέτησή τους.
  Αποτρέπουμε την καθολική σάρωση και τη συμπύκνωση των αντικειμένων του \selam SH\selg{} κατά τη συλλογή απορριμμάτων.
  Χρησιμοποιούμε έναν \selam card table\selg{} προσαρμοσμένο στα αποθηκευτικά μέσα για να παρακολουθούμε τις αναφορές από τον \selam SH\selg{} προς τον \selam DH\selg, περιορίζοντας τις σαρώσεις του συλλέκτη στα αντίστοιχα τμήματα του \selam SH\selg.
  Ανακτούμε αποθηκευτικό χώρο ανά περιοχή, αποδεχόμενοι αυξημένες απαιτήσεις χωρητικότητας για να περιορίσουμε τις λειτουργίες εισόδου/εξόδου που προκαλεί η συλλογή απορριμμάτων.
  Οι μηχανισμοί αυτοί μειώνουν συνολικά το κόστος σειριοποίησης, αποσειριοποίησης και συλλογής απορριμμάτων, βελτιώνοντας τη ρυθμαπόδοση των εφαρμογών και επεκτείνοντας τον διαχειριζόμενο σωρό πέρα από τη χωρητικότητα της \selam DRAM\selg.

  Στη συνέχεια, επεκτείνουμε τη χρήση σωρών σε αποθηκευτικά μέσα σε εφαρμογές που είναι ευαίσθητες στον χρόνο απόκρισης και βασίζονται σε συλλέκτες με σύντομες παύσεις.
  Επειδή η μεταφορά αντικειμένων στον \selam SH\selg{} μπορεί να αυξήσει τη διάρκεια των παύσεων συλλογής απορριμμάτων, επεκτείνουμε το μοντέλο επιλογής περιοχών του συλλέκτη, ώστε να συνυπολογίζει το κόστος μεταφοράς και τον χώρο που απελευθερώνεται στη \selam DRAM\selg{} τόσο από την ανάκτηση της μνήμης νεκρών αντικειμένων όσο και από τη μεταφορά ζωντανών αντικειμένων στον \selam SH\selg.
  Πραγματοποιούμε τη μεταφορά αντικειμένων μέσω παράλληλης εκκένωσης και παρέχουμε ξεχωριστό μηχανισμό για τη μεταφορά μεγάλων αντικειμένων που δεν εκκενώνονται κατά τις σταδιακές συλλογές.
  Οι μηχανισμοί αυτοί περιορίζουν την επιπλέον διάρκεια των παύσεων που οφείλεται στη μεταφορά αντικειμένων, μειώνοντας τον χρόνο απόκρισης των εφαρμογών στα υψηλά εκατοστημόρια (\selam tail latency\selg) και επεκτείνοντας τη χωρητικότητα του σωρού πέρα από τη χωρητικότητα της \selam DRAM\selg.

  Τέλος, εισάγουμε έναν μηχανισμό δυναμικής προσαρμογής του μεγέθους του σωρού που λαμβάνει υπόψη το κόστος εισόδου/εξόδου, με στόχο την αποδοτική κατανομή της \selam DRAM\selg{} μεταξύ του σωρού της \selam JVM\selg{} και της κρυφής μνήμης σελίδων.
  Ένας μεγαλύτερος σωρός μπορεί να μειώσει το κόστος συλλογής απορριμμάτων, αλλά αφήνει λιγότερη \selam DRAM\selg{} για την κρυφή μνήμη σελίδων και ενδέχεται να αυξήσει τον χρόνο αναμονής για είσοδο/έξοδο.
  Ο μηχανισμός εκτιμά το κόστος συλλογής απορριμμάτων και αναμονής για είσοδο/έξοδο και προσαρμόζει δυναμικά το μέγεθος του σωρού, ώστε να ελαχιστοποιήσει τη συνολική επιβάρυνση, διατηρώντας σταθερή τη συνολική διαθέσιμη \selam DRAM\selg.
  Η προσέγγιση αυτή θα μπορούσε να λειτουργήσει συμπληρωματικά προς τους σωρούς σε αποθηκευτικά μέσα, όπου ο \selam DH\selg{} ανταγωνίζεται για τη διαθέσιμη \selam DRAM\selg{} με την κρυφή μνήμη σελίδων που εξυπηρετεί τις προσβάσεις στον \selam SH\selg.

  Αξιολογούμε τους προτεινόμενους σχεδιασμούς σε εφαρμογές που δίνουν προτεραιότητα στη ρυθμαπόδοση και σε εφαρμογές ευαίσθητες στον χρόνο απόκρισης, όπως πλαίσια ανάλυσης δεδομένων, συστήματα επεξεργασίας γράφων, μηχανές αναζήτησης και συστήματα αποθήκευσης ζευγών κλειδιού-τιμής.
  Με την ίδια διαθέσιμη \selam DRAM\selg, ο σωρός σε αποθηκευτικό μέσο που στοχεύει στη ρυθμαπόδοση μειώνει τον χρόνο εκτέλεσης έως και κατά 73\% στο \selam Spark\selg{} και 28\% στο \selam Giraph\selg, σε σύγκριση με τις διαμορφώσεις τους που χρησιμοποιούν κρυφές μνήμες εκτός σωρού.
  Στις εφαρμογές που είναι ευαίσθητες στον χρόνο απόκρισης, ο προτεινόμενος σωρός σε αποθηκευτικό μέσο επιτυγχάνει έως και $161$ φορές μικρότερο χρόνο απόκρισης στο 99ο εκατοστημόριο (\selam P99\selg) σε σύγκριση με κρυφές μνήμες εντός σωρού που περιορίζονται από τη διαθέσιμη \selam DRAM\selg, $34$ φορές μικρότερο σε σύγκριση με κρυφές μνήμες εκτός σωρού που χρησιμοποιούν αποθηκευτικά μέσα και $9$ φορές μικρότερο σε σύγκριση με τον σχεδιασμό μας που στοχεύει στη ρυθμαπόδοση.
  Ο μηχανισμός προσαρμογής του μεγέθους του σωρού με βάση το κόστος εισόδου/εξόδου βελτιώνει την απόδοση των εφαρμογών κατά μέσο όρο κατά 30\% στο \selam Elasticsearch\selg{} και 33\% στο \selam Spark\selg, σε σύγκριση με τη βασική πολιτική του συλλέκτη απορριμμάτων για την προσαρμογή του μεγέθους του σωρού.
  Συνολικά, οι συνεισφορές αυτές δείχνουν πώς ο σχεδιασμός σωρών της \selam JVM\selg{} σε αποθηκευτικά μέσα μπορεί να βελτιώσει τη ρυθμαπόδοση και τον χρόνο απόκρισης στα υψηλά εκατοστημόρια σε εφαρμογές επεξεργασίας μεγάλου όγκου δεδομένων, ενώ η δυναμική προσαρμογή του μεγέθους του σωρού μπορεί να βελτιώσει την αξιοποίηση της \selam DRAM\selg.

  \paragraph{{\fontfamily{udidot}\selectfont Λέξεις-κλειδιά:}} εικονική μηχανή \selam Java (JVM)\selg, σωροί σε αποθηκευτικά μέσα, συλλογή απορριμμάτων, προσαρμογή μεγέθους σωρού με βάση το κόστος εισόδου/εξόδου, κρυφή μνήμη σελίδων του λειτουργικού συστήματος.

  \begin{center}
	  Επόπτης: Άγγελος Μπίλας \\ Καθηγητής \\ Τμήμα Επιστήμης Υπολογιστών \\ Πανεπιστήμιο Κρήτης
  \end{center}

 }

\selam


